%% Use the option review to obtain double line spacing
\documentclass[10pt]{article}
\usepackage[margin=1.2in]{geometry}
\usepackage{hyperref}
%\usepackage[demo]{graphicx}
\usepackage{graphicx}
\usepackage{subcaption}
\usepackage{cleveref}
\usepackage{comment}
%\usepackage{xr-hyper}
\newcommand{\beginsupplement}{
    \setcounter{section}{0}
    \renewcommand{\thesection}{S\arabic{section}}
    \setcounter{equation}{0}
    \renewcommand{\theequation}{S\arabic{equation}}
    \setcounter{table}{0}
    \renewcommand{\thetable}{S\arabic{table}}
    \setcounter{figure}{0}
    \renewcommand{\thefigure}{S\arabic{figure}}
    \newcounter{SIfig}
    \renewcommand{\theSIfig}{S\arabic{SIfig}}}


\begin{document}

%% Title, authors and addresses
%\title{Supplementary Information: \\ \\ Combining optimization and fire simulation modeling to protect biodiversity values at landscape scale}

%\author{Jaime Carrasco$^{a,b,}$\footnote{To whom correspondence should be addressed at Industrial Engineering Department, University of Chile, Santiago, Chile, E-mail: \url{jaimecarrasco@ug.uchile.cl}}, Rodrigo Mahaluf$^{a}$, Cristobal Pais$^{b,e}$, David Palacios$^{a}$, Fulgencio Lisón$^{c,d}$, Alejandro Miranda$^{d,f}$, and Andr\'es Weintraub$^{a,b}$ }

%\maketitle  
\begin{comment}
\noindent
{$^{a}$University of Chile, Industrial Engineering Department, Santiago, Chile \\
$^{b}$Complex Engineering System Institute - ISCI, Santiago, Chile. \\ 
$^{c}$Wildlife Ecology and Conservation Lab, Departamento de Zoología, Fac. Ciencias Naturales y Oceanográficas, Universidad de Concepción, Casilla 160-C, Concepción, Chile \\
$^{d}$Universidad de La Frontera, Departamento de Ciencias Forestales, Laboratorio de Ecología del Paisaje y Conservación, Temuco, Chile. \\
$^{e}$University of California Berkeley, IEOR Department, Berkeley, USA.} \\
$^{f}$University of Chile, Center for Climate and Resilience Research (CR$^{2}$), Santiago, Chile. \\
\end{comment}

%\end{comment}
%\clearpage

%%%%%%%%%%%%%%%%%%%%%%%%%% Instances

\beginsupplement
\noindent
%\textbf{Supplementary Information}
%\label{S:SuplementaryMaterial}
\noindent
\section*{Supplementary Information of: \\ \\ ``A two-stage stochastic programming approach incorporating spatially-explicit fire scenarios for optimal firebreak placement"}

Jaime Carrasco$^{a,b,}$\footnote{To whom correspondence should be addressed at Industrial Engineering Department, University of Chile, Santiago, Chile, E-mail: \url{jaimecarrasco@ug.uchile.cl}}, Rodrigo Mahaluf$^{a,b}$, Fulgencio Lisón$^{c,d}$, Cristobal Pais$^{b,e}$, Alejandro Miranda$^{d,f}$, Felipe de la Barra$^{a}$, David Palacios$^{a,b}$, and Andr\'es Weintraub$^{a,b}$ \\

\noindent
{$^{a}$University of Chile, Industrial Engineering Department, Santiago, Chile \\
$^{b}$Complex Engineering System Institute - ISCI, Santiago, Chile. \\ 
$^{c}$Wildlife Ecology and Conservation Lab, Departamento de Zoología, Fac. Ciencias Naturales y Oceanográficas, Universidad de Concepción, Casilla 160-C, Concepción, Chile \\
$^{d}$Universidad de La Frontera, Departamento de Ciencias Forestales, Laboratorio de Ecología del Paisaje y Conservación, Temuco, Chile. \\
$^{e}$University of California Berkeley, IEOR Department, Berkeley, USA.} \\
$^{f}$University of Chile, Center for Climate and Resilience Research (CR$^{2}$), Santiago, Chile. \\

\begin{figure*}[h!]
	%\centering
        \centerline{\includegraphics[scale=1]{Pictures/violin_plot.pdf}}
        \caption{Distribution of the values of wind speed and direction, temperature and relative humidity that are included in the 100 fire weather scenarios that were used in our study.}
	\label{fig:violin_plot}
\end{figure*}

\begin{figure*}[h!]
	%\centering
    \centerline{\includegraphics[scale=.54]{Pictures/Combined_index.pdf}}
    \caption{Combined Index Map.}
	\label{fig:IdexResults}
\end{figure*}

\begin{figure*}[h!]
	%\centering
    \centerline{\includegraphics[scale=.54]{Pictures/IPM.pdf}}
    \caption{Fire Ignition Probability Map.}
	\label{fig:IPM}
\end{figure*}

\begin{figure*}[h!]
	%\centering
    \centerline{\includegraphics[scale=.54]{figs/LN_DPV.pdf}}
    \caption{Map of the Downstream Protection Value natural logarithm.}
	\label{fig:DPV}
\end{figure*}

\begin{figure*}[h!]
	%\centering
    \centerline{\includegraphics[scale=.54]{figs/DPV_zoom.pdf}}
    \caption{Downstream Protection Value Map.}
	\label{fig:log_DPV}
\end{figure*}



\begin{table}[!ht]
\centering
\caption{Number of bird's occurrences by species, sorted in decreasing order.}
\resizebox{\textwidth}{!}{
\begin{tabular}{rlr | rlr}
  \hline
\# & Especie & Obs. & \# & Especie & Obs.\\ 
  \hline
  1 & \textit{Milvago chimango} & 265 & 29 & \textit{Podilymbus podiceps} &  87 \\ 
    2 & \textit{Theristicus melanopis} & 250 & 30 & \textit{Larus dominicanus} &  86\\
    3 & \textit{Vanellus chilensis} & 243 & 31 & \textit{Anas flavirostris} &  79 \\ 
    4 & \textit{Elaenia albiceps} & 188 & 32 & \textit{Enicognathus leptorhynchus} &  79\\
    5 & \textit{Aphrastura spinicauda} & 184 & 33 & \textit{Notiochelidon cyanoleuca} &  79 \\  
    6 & \textit{Tachycineta leucopyga} & 174 &  34 & \textit{Passer domesticus} &  79 \\  
    7 & \textit{Turdus falcklandii} & 164 & 35 & \textit{Scytalopus magellanicus} &  79\\ 
    8 & \textit{Sephanoides sephanoides} & 146 & 36 & \textit{Mimus thenca} &  78 \\
    9 & \textit{Troglodytes aedon} & 145 & 37 & \textit{Zonotrichia capensis} &  75\\
   10 & \textit{Scelorchilus rubecula} & 143 & 38 & \textit{Megaceryle torquata} &  68 \\
   11 & \textit{Anairetes parulus} & 127 & 39 & \textit{Campephilus magellanicus} &  62 \\
   12 & \textit{Pygarrhichas albogularis} & 119 & 40 & \textit{Anas georgica} &  60 \\
   13 & \textit{Spinus barbatus} & 115 & 41 & \textit{Ardea alba} &  55 \\  
   14 & \textit{Cinclodes patagonicus} & 112 & 42 & \textit{Sturnella loyca} &  55 \\
   15 & \textit{Phrygilus patagonicus} & 110 & 43 & \textit{Columba livia} &  54 \\
   16 & \textit{Fulica armillata} & 109 & 44 & \textit{Zenaida auriculata} &  52 \\
   17 & \textit{Podiceps major} & 108 & 45 & \textit{Buteo polyosoma} &  50  \\
   18 & \textit{Chroicocephalus maculipennis} & 106 & 46 & \textit{Colorhamphus parvirostris} &  47 \\
   19 & \textit{Xolmis pyrope} & 106 & 47 & \textit{Sicalis luteola} &  40 \\ 
   20 & \textit{Patagioenas araucana} & 103 & 48 & \textit{Cinclodes fuscus} &  37 \\
   21 & \textit{Phalacrocorax brasilianus} &  96 & 49 & \textit{Sylviorthorhynchus desmursii} &  36 \\
   22 & \textit{Pteroptochos tarnii} &  94 & 50 & \textit{Cathartes aura} &  34 \\
   23 & \textit{Coragyps atratus} &  93 & 51 & \textit{Dryobates lignarius} &  34 \\
   24 & \textit{Curaeus curaeus} &  93 & 52 & \textit{Glaucidium nana} &  33 \\
   25 & \textit{Egretta thula} &  93 & 53 & \textit{Chloephaga poliocephala} &  30 \\ 
   26 & \textit{Enicognathus ferrugineus} &  93 & 54 & \textit{Falco sparverius} &  30 \\
   27 & \textit{Colaptes pitius} &  92 & 55 & \textit{Molothrus bonariensis} &  30 \\
   28 & \textit{Caracara plancus} &  87 \\  
   \hline
\end{tabular}}
\label{tab:birds}
\end{table}


\end{document}